% !TEX TS-program = lualatex
% !TEX encoding = UTF-8 Unicode
\documentclass[ngerman,11pt]{article}
\usepackage{luacode}
\usepackage{fontspec}
\usepackage{emoji}
\usepackage[ngerman]{babel}
\usepackage[ngerman]{datetime2}
\usepackage[a4paper, total={16cm, 25cm}]{geometry}
\usepackage{amsmath}
\usepackage{amsfonts}
\usepackage{amssymb}
\usepackage{multicol}
\usepackage{tcolorbox}
\usepackage{cancel}
\usepackage{enumitem}
\usepackage{hyperref}
\usepackage{array}
\usepackage{xcolor}
\usepackage{listings}

\definecolor{codebg}{HTML}{F4F6F8}
\definecolor{codecomment}{HTML}{6A737D}
\definecolor{codekeyword}{HTML}{0550AE}
\definecolor{codestring}{HTML}{0A7D5C}

\lstdefinestyle{python}{
    language=Python,
    basicstyle=\ttfamily\footnotesize,
    backgroundcolor=\color{codebg},
    commentstyle=\color{codecomment}\itshape,
    keywordstyle=\color{codekeyword}\bfseries,
    stringstyle=\color{codestring},
    showstringspaces=false,
    frame=single,
    rulecolor=\color{codebg},
    xleftmargin=6pt,
    xrightmargin=6pt,
    aboveskip=6pt,
    belowskip=6pt,
    breaklines=true,
    numbers=none,
}

\begin{document}
\hrule
\begin{center}
{\large\bf Das Perzeptron-Lernverfahren}\\[-0mm]
Studienkolleg der TU Berlin\hfill Till Zoppke\\
Informatik, Oberkurs \hfill \DTMgermanmonthname{\month} \the\year\\[2mm]
\end{center}
\hrule
\bigskip

\noindent Ein \textbf{Perzeptron} (Rosenblatt, 1958) ist das einfachste künstliche Neuron, das aus Beispielen lernen kann. Es berechnet aus den Eingaben $x_1, x_2$ eine gewichtete Summe $z = w_0 + w_1 x_1 + w_2 x_2$ und gibt als Vorhersage das Vorzeichen von $z$ zurück. Das Gewicht $w_0$ fungiert dabei als \emph{Bias}. Während des Trainings werden die Gewichte bei jeder Fehlklassifikation angepasst, bis alle Trainingsbeispiele korrekt klassifiziert sind.

\begin{tcolorbox}[title=Der Perzeptron-Lernalgorithmus (Pseudocode)]
\textbf{Trainingsdaten:} Beispiele $(x_1^{(k)}, x_2^{(k)}, y^{(k)})$ für $k = 1, \ldots, m$ mit $y^{(k)} \in \{-1, +1\}$

\medskip
\textbf{Initialisierung:} $w_0 = 0$, $w_1 = 0$, $w_2 = 0$

\medskip
\textbf{Wiederhole} für jede \emph{Epoche}, bis keine Fehler mehr auftreten:
\begin{quote}
\textbf{Für jedes} Trainingsbeispiel $(x_1, x_2, y)$:
\begin{quote}
Berechne $z = w_0 + w_1 x_1 + w_2 x_2$ und $\hat{y} = \text{sign}(z)$\\
\textbf{Falls} $\hat{y} \neq y$ \quad (Fehlklassifikation):
\begin{quote}
$w_0 \leftarrow w_0 + y$; \quad $w_1 \leftarrow w_1 + y \cdot x_1$; \quad $w_2 \leftarrow w_2 + y \cdot x_2$
\end{quote}
\end{quote}
\end{quote}

\medskip
\textbf{Ausgabe:} Gewichte $w_0, w_1, w_2$, die alle Trainingsbeispiele korrekt klassifizieren.
\end{tcolorbox}

\subsection*{Aufgabe 1: Algorithmus per Hand anwenden (OR-Funktion)}

\noindent Trainieren Sie ein Perzeptron so, dass es die OR-Funktion lernt. Verwenden Sie die Codierung $\{-1, +1\}$.

\medskip
\noindent \textbf{a)} Stellen Sie zunächst die Trainingsdaten für die OR-Funktion als Tabelle auf.

\medskip
\noindent \textbf{b)} Führen Sie den Algorithmus per Hand durch. Tragen Sie in jeder Zeile die Nettoeingabe $z$ und die neuen Gewichte nach einem eventuellen Update ein. Notieren Sie bei jedem Update explizit, \emph{wie} die Gewichte angepasst werden.

\medskip
\noindent\begin{tabular}{@{}p{1.8cm}@{}l@{}}
Epoche 1 &
\begin{tabular}{|>{\centering\arraybackslash}p{2.2cm}|>{\centering\arraybackslash}p{1.2cm}|>{\centering\arraybackslash}p{1.5cm}|p{7cm}|}
\hline
$(x_1, x_2)$ & $y$ & $z$ & Update (falls nötig) und neue Gewichte \\
\hline
Start & -- & -- & $w_0 = 0, \; w_1 = 0, \; w_2 = 0$ \\
\hline
 & & & \\[8pt]
\hline
 & & & \\[8pt]
\hline
 & & & \\[8pt]
\hline
 & & & \\[8pt]
\hline
\end{tabular}
\end{tabular}

\medskip
\noindent\begin{tabular}{@{}p{1.8cm}@{}l@{}}
Epoche 2 &
\begin{tabular}{|>{\centering\arraybackslash}p{2.2cm}|>{\centering\arraybackslash}p{1.2cm}|>{\centering\arraybackslash}p{1.5cm}|p{7cm}|}
\hline
$(x_1, x_2)$ & $y$ & $z$ & Update (falls nötig) und neue Gewichte \\
\hline
 & & & \\[8pt]
\hline
 & & & \\[8pt]
\hline
 & & & \\[8pt]
\hline
 & & & \\[8pt]
\hline
\end{tabular}
\end{tabular}

\newpage

% ===========================================================
%                      SEITE 2: GLOSSAR
% ===========================================================

\section*{Glossar: Perzeptron-Lernverfahren}

\begin{center}
\begin{tabular}{|>{\raggedright\arraybackslash}p{4.2cm}|>{\raggedright\arraybackslash}p{10.8cm}|}
\hline
\textbf{Begriff} & \textbf{Erklärung} \\
\hline
Aktivierungsfunktion & Funktion, die aus der Nettoeingabe $z$ die Ausgabe erzeugt. Beim Perzeptron: die Vorzeichenfunktion $\text{sign}(z)$ mit Werten $+1$ oder $-1$. \\
\hline
Bias & Gewicht $w_0$ ohne zugehörige Eingabe. Verschiebt die Entscheidungsgrenze vom Ursprung weg. Ohne Bias ginge die Trenngerade zwangsläufig durch den Ursprung. \\
\hline
Epoche & Ein vollständiger Durchlauf durch alle Trainingsbeispiele. Das Training besteht meist aus mehreren Epochen. \\
\hline
Fehlklassifikation & Tritt auf, wenn $\hat{y} \neq y$, also wenn Vorhersage und wahre Klasse nicht übereinstimmen. \\
\hline
Gewicht (\emph{weight}) & Parameter $w_i$, der angibt, wie stark die Eingabe $x_i$ in die Nettoeingabe eingeht. Gewichte werden durch das Training angepasst. \\
\hline
Klassifikationsproblem & Aufgabe, Eingaben einer von mehreren vorgegebenen Klassen zuzuordnen. Bei einem \emph{binären} Klassifikationsproblem gibt es genau zwei Klassen, hier codiert als $+1$ und $-1$. Das Perzeptron löst ausschließlich binäre Klassifikationsprobleme. \\
\hline
Nettoeingabe & Gewichtete Summe $z = w_0 + w_1 x_1 + w_2 x_2$. Wert \emph{vor} der Aktivierungsfunktion. \\
\hline
Neuron & Grundeinheit eines künstlichen neuronalen Netzes. Berechnet eine gewichtete Summe der Eingaben und wendet darauf eine Aktivierungsfunktion an. \\
\hline
Training & Der Vorgang, bei dem die Gewichte $w_0, w_1, w_2$ schrittweise angepasst werden, bis das Neuron alle Trainingsdaten korrekt klassifiziert. \\
\hline
Trainingsdaten & Die Menge aller Beispiele $(x_1, x_2, y)$, mit denen das Neuron lernt. Jeder Datensatz besteht aus den Eingabewerten und dem zugehörigen Zielwert $y$. \\
\hline
Vorhersage & Ergebnis der Aktivierungsfunktion: $\hat{y} = \text{sign}(z) \in \{-1, +1\}$. Der Hut über $y$ kennzeichnet den vorhergesagten Wert im Unterschied zum wahren Zielwert $y$. \\
\hline
Überwachtes Lernen & Lernverfahren, bei dem jedes Trainingsbeispiel mit einem bekannten Zielwert $y$ versehen ist. Das Modell lernt, indem es seine Vorhersagen $\hat{y}$ mit den wahren Werten vergleicht und bei Fehlern die Gewichte anpasst. Das Perzeptron ist ein klassisches Beispiel für überwachtes Lernen. \\
\hline
\end{tabular}
\end{center}

\bigskip
\subsection*{Zum Ausprobieren}

\begin{itemize}[leftmargin=*]
\item \textbf{Perzeptron-Simulator (Universität Passau):} Visualisiert den Lernalgorithmus Schritt für Schritt. Gewichtsvektor und Trenngerade werden live dargestellt.\\
\url{https://kex.fim.uni-passau.de/Perzeptron_Online/html/Perzeptron_Online.html}

\item \textbf{TensorFlow Playground:} Interaktiver Simulator für mehrschichtige Netze. Zum Experimentieren mit nicht-linear separierbaren Datensätzen (z.\,B.\ XOR-artige Spiralen).\\
\url{https://playground.tensorflow.org/}
\end{itemize}

\bigskip
\noindent\footnotesize Lizenz: CC BY-SA 4.0 $\cdot$ Studienkolleg TU Berlin $\cdot$ Till Zoppke

\end{document}
